% called by main.tex
% Glosario en preliminares (movido desde los anexos el 2026-07-10, punto 1 de la tutora)
\section*{Glosario de términos}
\addcontentsline{toc}{section}{Glosario de términos}
\label{sec::glosario}

Glosario pensado para el lector no familiarizado con las criptomonedas ni con la
microestructura de mercado. Recoge los términos empleados a lo largo de la memoria.

\begin{description}[align=left, leftmargin=3.2cm, labelwidth=3cm, font=\bfseries]
\item[Mercado de predicción] Mercado en el que se negocian contratos cuyo valor depende del
resultado de un acontecimiento futuro.
\item[Contrato binario] Contrato que paga una unidad si el suceso ocurre y cero si no; su
precio toma valores entre 0 y 1 y se interpreta como probabilidad.
\item[Polymarket] Plataforma de mercados de predicción descentralizada, con emparejamiento
fuera de cadena y liquidación en cadena.
\item[Libro de órdenes (LOB)] Lista, ordenada por precio, de las órdenes de compra y venta
pendientes de ejecución.
\item[Bid / Ask] Mejor precio de compra (\emph{bid}) y mejor precio de venta (\emph{ask}).
\item[Spread] Diferencia entre \emph{ask} y \emph{bid}; primer coste de cruzar el mercado.
\item[Profundidad] Volumen disponible en cada nivel de precio del libro.
\item[Precio medio (\emph{mid})] Punto medio entre \emph{bid} y \emph{ask}.
\item[Tick] Variación mínima de precio del mercado; unidad de medida de los resultados.
\item[Microestructura] Estudio de cómo se forman los precios a partir de órdenes, \emph{spread},
profundidad y flujo, a escalas de tiempo muy finas.
\item[Markout] Movimiento del precio desde el instante de ejecución a lo largo de un horizonte;
sirve para medir si una operación fue buena.
\item[Latencia] Retardo entre observar la oportunidad y ejecutar la orden.
\item[Spot / Perpetuo] Mercado de Bitcoin al contado (\emph{spot}) y de futuros sin vencimiento
(\emph{perpetual futures}); referencias externas del contrato.
\item[Oráculo de precios] Fuente externa que aporta el precio de referencia de Bitcoin.
\item[Volatilidad realizada] Medida de cuánto se ha movido el precio recientemente.
\item[Ventana \emph{prestart}] Intervalo inmediatamente anterior a la apertura del mercado.
\item[Validación temporal] Partición de los datos respetando el orden del tiempo (entrenar con
el pasado, evaluar con el futuro).
\item[Fuera de muestra (OOS)] Datos no usados ni para entrenar ni para seleccionar.
\item[Fuera de pliegue (OOF)] Protocolo en el que las predicciones que consume una segunda etapa
proceden siempre de un modelo que no las entrenó.
\item[Fuga temporal] Filtración de información del futuro hacia el entrenamiento o las
características; principal fuente de resultados ilusorios.
\item[Congelación] Fijar el modelo y los umbrales antes de observar el conjunto fuera de muestra.
\item[Drawdown] Mayor caída acumulada desde un máximo previo.
\item[\emph{Bootstrap}] Remuestreo con reemplazo para estimar la incertidumbre de un resultado.
\item[Registro de sombra] Contabilidad de lo que la política habría indicado sin arriesgar
capital. Puede ser retrospectivo (replay \emph{offline}) o prospectivo; solo el segundo aporta
nuevos días de validación.
\item[Gradient boosting] Familia de modelos que suma muchos árboles de decisión, cada uno
corrigiendo los errores del anterior.
\item[Taker / Maker] Quien cruza el libro y ejecuta contra órdenes existentes (\emph{taker},
paga la comisión) frente a quien aporta liquidez con órdenes limitadas (\emph{maker}, exento
de comisión y receptor de \emph{rebates}).
\item[Selección adversa] Riesgo del \emph{maker} de ser ejecutado precisamente cuando el
precio va a moverse en su contra.
\item[Riesgo terminal] En un contrato binario, el riesgo del inventario que llega vivo a
la resolución: el contrato liquida a 0 o a 1, de modo que una posición sin cerrar se
enfrenta a la varianza estructural del desenlace.
\item[Modelo fundacional] Modelo pre-entrenado sobre colecciones masivas y heterogéneas
de datos (series temporales, tablas) pensado para aplicarse a tareas nuevas sin
entrenamiento o con una adaptación ligera.
\item[Zero-shot / linear-probing / fine-tuning] Regímenes de adaptación crecientes de un
modelo pre-entrenado: usarlo tal cual, entrenar solo su cabeza de salida, o re-entrenar
todos sus parámetros.
\item[Puntuación de Brier] Métrica de calidad de una probabilidad (\emph{proper scoring}):
error cuadrático medio entre la probabilidad emitida y el desenlace; 0,25 equivale a una
moneda al aire.
\item[Pre-registro] Compromiso público de hipótesis, protocolo y criterios de éxito
anterior a la observación de los datos de evaluación.
\item[Sello de tiempo (OpenTimestamps)] Atestación criptográfica sobre la cadena de
Bitcoin de que un documento existía antes de cierto bloque; permite verificar la
anterioridad de un compromiso sin confiar en el autor.
\end{description}
